\documentclass{amsproc}
\usepackage{mathptmx}      % use Times fonts if available on your TeX system
\usepackage{latexsym}
\usepackage{amssymb}
\usepackage{bm}
\usepackage[usenames,dvipsnames]{xcolor}
\usepackage{graphicx,pifont}
\usepackage{multicol,multirow}
\usepackage{ulem,cancel}
\usepackage{tikz}
\usetikzlibrary[patterns]
\usepackage{comment}
\usepackage{hyperref}
\usepackage{array}
\usepackage{caption}
\usepackage{subcaption}
\usepackage{txfonts}
\usepackage{enumerate}
\usepackage{listings}

\newtheorem{theorem}{Theorem}[section]
\newtheorem{lemma}[theorem]{Lemma}
\newtheorem{proposition}{Proposition}[section]
\theoremstyle{definition}
\newtheorem{definition}[theorem]{Definition}
\newtheorem{example}[theorem]{Example}
\theoremstyle{remark}
\newtheorem{remark}[theorem]{Remark}
\numberwithin{equation}{section}
\numberwithin{figure}{section}
\newtheorem{conjecture}[theorem]{Conjecture}

\newcommand{\crossref}{{\color{red}ref here}}

\newcommand{\bs}[1]{\boldsymbol{#1}}
\newcommand{\bsb}[1]{{\mbox{\sffamily #1}}}
\newcommand{\bb}[1]{{\mbox{\sffamily \bfseries #1}}}
\newcommand{\oname}[1]{\textrm{#1}}
\newcommand{\abs}[1]{\left|#1\right|}
\newcommand{\nrm}[1]{\left|\left|#1\right|\right|}
\newcommand{\eqdef}{\stackrel{\mathrm{def}}{=\joinrel=}}
\newcommand{\ab}[1]{\langle#1\rangle}
\newcommand{\po}[1]{\frac{\partial}{\partial{#1}}}
\newcommand{\poh}[1]{\partial/\partial{#1}}
\newcommand{\pp}[2]{\frac{\partial {#1}}{\partial {#2}}}
\newcommand{\pph}[2]{\partial {#1}/\partial {#2}}
\newcommand{\termc}[1]{{\mbox{\sffamily {#1}}}}
\newcommand{\bt}[1]{{\boldsymbol{\mbox{\sffamily \bfseries {#1}}}}}
\newcommand{\phf}[1]{{#1}+1/2}
\newcommand{\mhf}[1]{{#1}-1/2}
\newcommand{\powth}{{\textrm{\scriptsize th}}}
\newcommand{\powrd}{{\textrm{\scriptsize rd}}}
\newcommand{\powst}{{\textrm{\scriptsize st}}}
\newcommand{\pownd}{{\textrm{\scriptsize nd}}}

\setlength{\unitlength}{1cm}
\setlength{\arraycolsep}{3pt} 
\setlength{\tabcolsep}{1pt}

\begin{document}

\title[\texttt{TumorDyn} Manual]{\bf A Brief Manual for \texttt{TumorDyn}}
\author[X.~Zeng]{Xianyi Zeng}
\address{Department of Mathematical Sciences, University of Texas at El Paso, El Paso TX 79902, United States.\\
         Tel.: +1-915-747-6759}
\email[X.~Zeng]{xzeng@utep.edu}
\date{\today}

%\subjclass[2010]{65M12 \and 35L65}

%\keywords{
%}

\begin{abstract}
This is a short guide for obtaining and using \texttt{TumorDyn}, a finite-volume solver for simulating spherically symmetric tumor growth problems.
It also contains a simple example to run several runs in parallel on a \texttt{macOS}/\texttt{Linux} machine with multiple processors, as well as some post-processing scripts using \texttt{Matlab}.
\end{abstract}

\maketitle

\setcounter{tocdepth}{2}
\tableofcontents

\section{Obtaining, Updating, and Compiling \texttt{TumorDyn}}

\texttt{TumorDyn} is written in \texttt{C++}; the project is housed at \texttt{Bitbucket} with the name \texttt{TDSolve}.
Everything in this manual assumes operation in terminal with proper shell commands, unless otherwise specified.
On latest versions of \texttt{macOS}, this needs to be enabled by installing ``Command Line Tools''; any \texttt{Linux} system should come with a shell command set already.
In order to get the source code from \texttt{Bitbucket} and compile it successfully, the following tools are required:
\begin{itemize}
\item \texttt{hg}: Mercurial, a distributed source control management tool used to obtaining and updating the source code from \texttt{Bitbucket} repository.
\item \texttt{make}: GNU Make, a tool to control the generation of executables and other non-source files from source codes.
\item \texttt{g++}: the GNU \texttt{C++} compiler. This can be replaced by \texttt{clang++} on \texttt{macOS} -- need to change a line in the \texttt{makefile} shipped with the package.
\item \texttt{cmake} (optional): New BSD CMake.
\end{itemize}
On \texttt{Linux}, these tools should be available through any package management tools such as \texttt{apt-get} on \texttt{Ubuntu} systems or \texttt{yum} on \texttt{Red Hat} based systems such as \texttt{Red Hat Linux}, \texttt{Fedora Linux}, \texttt{CentOS}, and \texttt{Scientific Linux}; on \texttt{macOS}, however, \texttt{MacPorts} is a good source to get many useful tools. 

\subsection{To Start: Getting \texttt{TumorDyn} from \texttt{Bitbucket} and Make!}
If you have not received an invitation from me or co-developers of the program, stop here, you won't be able to obtain the code. 
Now we assume you already have received an invitation, joined the project, and registered an \texttt{Atlassian} account.
To begin with, let us create an empty folder, such as \texttt{\$HOME/td\_proj}; anyway let us call this folder \texttt{\$TDHOME}.
While in this folder and with \texttt{hg} installed, use the following to create a local copy of the program:
\begin{lstlisting}[backgroundcolor = \color{lightgray}, language = bash, xleftmargin = 0.1cm, framexleftmargin = 0.0cm]
$ hg clone ssh://hg@bitbucket.org/xyzeng/tdsolve TDSolve
\end{lstlisting}
Here and throughout the manual, such a line beginning with \$ on a shaded background indicates a command to type in the terminal; without the \$ sign, it is just a line showing some output on the terminal, such as by listing the content of the current directory.
This line in particular, will create a folder called \texttt{TDSolve} in our project home \texttt{\$TDHOME}; feel free to use any other name you like to replace \texttt{TDSolve}.
Note that mercurial use \texttt{ssh} to connect to the \texttt{Bitbucket} repository; if this is your first time doing so you may get a few error messages.
When this happens, please refer to the \texttt{Atlassian} documentation page for resolving the issues:
\begin{center}
{\color{blue} \url{https://confluence.atlassian.com/bitbucket/troubleshoot-ssh-issues-271943403.html}}
\end{center}
The \texttt{hg clone} command may take a while, as you may have expected, since it copies everything from the repository to the local machine.
Once it is done, enter this directory, you will see a few files and sub-folders; among other things, there are:
\begin{lstlisting}[backgroundcolor = \color{lightgray}, language = bash, xleftmargin = 0.1cm, framexleftmargin = 0.1cm]
include/ main/ makefile readme.txt src/ tools/
CMakeLists.txt build_cmake.sh
\end{lstlisting}
The file \texttt{readme.txt} contains an extremely short guide to use the program, assuming that you're already familiar with the bash commands, etc.
The three first three folders contain the source codes: \texttt{include/} contains all the header files, \texttt{src/} contains all the source files excepting the one containing the \texttt{main} function, which is in the folder \texttt{main/}.
You may notice that there is also a folder called \texttt{obj/}, this is used to hold all the object files generated during compiling the program; no need to worry about it.
The folder \texttt{tools/} contains some \texttt{Matlab} scripts that are used to read the simulation results and do some simple plots; we will discuss this later.

There are two options to build the program.
\begin{enumerate}[Option 1]
  \item If you have \texttt{cmake}, type:
\begin{lstlisting}[backgroundcolor = \color{lightgray}, language = bash, xleftmargin = 0.1cm, framexleftmargin = 0.0cm]
$ sh build_cmake.sh
\end{lstlisting}
  while in the folder \texttt{\$TDHOME/TDSolve}; it will automatically search for the default \texttt{C++} compiler in the system and build the program.
  \item If you do not have \texttt{cmake} or want to specify the compiler, the \texttt{makefile} should be used.
  To do so, type:
  \begin{lstlisting}[backgroundcolor = \color{lightgray}, language = bash, xleftmargin = 0.1cm, framexleftmargin = 0.0cm]
  $ make 
  \end{lstlisting}
  while in the folder \texttt{\$TDHOME/TDSolve}; it will automatically call the \texttt{C++} compiler and build the program.
  By default, the GNU compiler \texttt{g++} is used; if you're using \texttt{macOS} and decide not to use \texttt{g++}, another option is to use \texttt{clang++} that is shipped with the system.
  To do so, uncomment the second line in \texttt{makefile} by deleting the first \# character, and comment the first line by adding a \# character at the beginning of the line.
\end{enumerate}

If there are more than one processors on the machine, the building process can be speed up a little bit by making use of these processors; to do so, for example with 4 processors, type instead:
\begin{lstlisting}[backgroundcolor = \color{lightgray}, language = bash, xleftmargin = 0.1cm, framexleftmargin = 0.0cm]
$ make -j 4
\end{lstlisting}
for the second option; and modify the script \texttt{build\_cmake.sh} accordingly for the first option.
Parallel building may actually be an overkill; anyway this is not a very large program, at least for now...

Upon a successful building, an executable file with the name:
\begin{lstlisting}[backgroundcolor = \color{lightgray}, language = bash, xleftmargin = 0.1cm, framexleftmargin = 0.1cm]
TumorDynRun
\end{lstlisting}
should be added to the current folder; this is the executable to run all the simulations.

\subsection{Keeping the Project Updated}
Occasionally the \texttt{Bitbucket} repository will be updated to include new features or to get rid of some obsolete features.
To keep the local copy updated, type:
\begin{lstlisting}[backgroundcolor = \color{lightgray}, language = bash, xleftmargin = 0.1cm, framexleftmargin = 0.0cm]
$ hg pull
\end{lstlisting}
in \texttt{\$TDHOME/TDSolve} to fetch all the updates.
If there are any conflicts reported by mercurial, these conflicts must be resolved and an \texttt{hg merge} is required before continuing; this should not happen if you do not modify any file in this folder except for those generated by \texttt{make}.
You can see if anything is changed by typing:
\begin{lstlisting}[backgroundcolor = \color{lightgray}, language = bash, xleftmargin = 0.1cm, framexleftmargin = 0.0cm]
$ hg st
\end{lstlisting}
in the current directory.
Once the \texttt{hg pull} is successful, use:
\begin{lstlisting}[backgroundcolor = \color{lightgray}, language = bash, xleftmargin = 0.1cm, framexleftmargin = 0.0cm]
$ hg up default
\end{lstlisting}
to update all the source files to the newest version (we're not worrying about any branch other than the default one for now).
Next, use:
\begin{lstlisting}[backgroundcolor = \color{lightgray}, language = bash, xleftmargin = 0.1cm, framexleftmargin = 0.0cm]
$ sh build_cmake.sh
\end{lstlisting}
or
\begin{lstlisting}[backgroundcolor = \color{lightgray}, language = bash, xleftmargin = 0.1cm, framexleftmargin = 0.0cm]
$ make 
\end{lstlisting}
or
\begin{lstlisting}[backgroundcolor = \color{lightgray}, language = bash, xleftmargin = 0.1cm, framexleftmargin = 0.0cm]
$ make -j 4
\end{lstlisting}
to re-build the code.
This process could be faster than the first build because GNU Make will automatically determine the minimum amount of work to re-compile the code -- that is, if a source file is not changed during the update the corresponding object (\texttt{.o}) file will not be re-produced.

As a final remark, if the folder \texttt{\$TDHOME/TDSolve} is not in the \texttt{PATH} variable, you'll need to provide the full path to \texttt{TumorDynRun} in order to use it.

\section{Input File Explained and A Simple Run}
Along with the source \texttt{C++} files, there is also a folder:
\begin{lstlisting}[backgroundcolor = \color{lightgray}, language = bash, xleftmargin = 0.1cm, framexleftmargin = 0.1cm]
manual/
\end{lstlisting}
in the cloned \texttt{\$TDHOME/TDSolve}.
This directory contains this manual as well as its \LaTeX files in the sub-folder \texttt{doc/}.
The other sub-folder, \texttt{demo/}, holds a \texttt{hcl} file called the input file.
Why the extension \texttt{.hcl}? Because many of the \texttt{C++} source files here are borrowed from a much larger project called \texttt{HCLSolve} (Hyperbolic Conservation Law Solvers)... So we just inhere the naming convention.

Preparing ourselves for a first run of \texttt{TumorDynRun}, let us go to \texttt{\$TDHOME} and create a new folder called \texttt{TDRun}:
\begin{lstlisting}[backgroundcolor = \color{lightgray}, language = bash, xleftmargin = 0.1cm, framexleftmargin = 0.0cm]
$ cd $TDHOME
$ mkdir TDRun
\end{lstlisting}
Of course, you can use any other name or even just run in \texttt{\$TDHOME/TDSolve/manual/demo}, where the input files exist;
however, I suggest not to put the simulation files in \texttt{\$TDHOME/TDSolve}, since this will create a lot of new files when you check the status of the program using \texttt{hg st}, unless we also modify the hidden file \texttt{\$TDHOME/TDSolve/.hgignore}...

\subsection{A Simple Run} 
Go to \texttt{\$TDHOME/TDRun} and create a new folder \texttt{TestRun}, enter it, and copy the file \texttt{input\_sample.hcl} here:
\begin{lstlisting}[backgroundcolor = \color{lightgray}, language = bash, xleftmargin = 0.1cm, framexleftmargin = 0.0cm]
$ cd $TDHOME/TDRun
$ mkdir TestRun 
$ cd TestRun
$ cp $TDHOME/TDSolve/manual/demo/input_sample.hcl ./
$ mv input_sample.hcl input.hcl
\end{lstlisting}
We have renamed it for simplicity.
Use the following line to run the program with \texttt{input.hcl} as the input file:
\begin{lstlisting}[backgroundcolor = \color{lightgray}, language = bash, xleftmargin = 0.1cm, framexleftmargin = 0.0cm]
$ $TDHOME/TDSolve/TumorDynRun -i input.hcl
\end{lstlisting}
Running this command will produce a program logo and some diagonostic information:
{\tiny
\begin{lstlisting}[backgroundcolor = \color{lightgray}, language = bash, xleftmargin = 0.1cm, framexleftmargin = 0.1cm]
Iter   0: Time 0.00000000; Time step 0.00000100; Radius 0.20000000; Front-velocity 0.05105135.
Iter   1: Time 0.00000100; Time step 0.00000120; Radius 0.20000005; Front-velocity 0.05194606.
Iter   2: Time 0.00000220; Time step 0.00000144; Radius 0.20000011; Front-velocity 0.05273993.
\end{lstlisting}
}
Each line contains from the left to the right:
\begin{itemize}
\item The iteration number $k$ starting with 0, the initial condition.
\item Current simulation time $t^k$.
\item The next time step size $\Delta t^k = t^{k+1}-t^k$.
\item Current radius of the tumor $R^k$.
\item Current growth velocity at the tumor boundary $V^k(R^k,t^k)$.
\end{itemize}
Once the run is completed (it will take a while), two solution files are generated:
\begin{lstlisting}[backgroundcolor = \color{lightgray}, language = bash, xleftmargin = 0.1cm, framexleftmargin = 0.1cm]
sol.dat rad.dat
\end{lstlisting}
They will be explained later.

\subsection{The Input File}
There are four sections of the input file we just used, they are separated by an empty line.
Each section is composed of a few lines that look like:
\begin{lstlisting}[backgroundcolor = \color{lightgray}, language = bash, xleftmargin = 0.1cm, framexleftmargin = 0.1cm]
// Some comments
\end{lstlisting}
or
\begin{lstlisting}[backgroundcolor = \color{lightgray}, language = bash, xleftmargin = 0.1cm, framexleftmargin = 0.1cm]
variableName = variableValue
\end{lstlisting}
The first scenario is the comments and it has nothing to do with the simulation. 
It is the second scenario that set various values/options in the simulation; it doesn't matter how many spaces there are between the equality sign and either side of the equation.

The first section assign values to the model parameters and they're fairly self-explanatory.
For example:
\begin{itemize}
\item \texttt{prob\_id} identifies the problem to be solved:
  \begin{itemize}
    \item 0: The PDGF model.
    \item 1: The OPN model.
  \end{itemize}
  This choice affects how the medel parameters are defined and the output variables.
  More details are offered in Section~\ref{sec:model}.
\item \texttt{n\_bc} is the boundary value $N_{\textrm{bc}}$ for the immune cell number $N$ (or $M$ of the OPN model).
%\item \texttt{n\_star} is the threshold value $N^*$ that determines the dissipation rate of $N$; these rates are given by \texttt{rho\_mnus} if during the computation $N$ is smaller than $N^*$, and by \texttt{rho\_plus} if $N$ is larger than $N^*$.
\item \texttt{logistic} is a switch for the logistic terms in the OPN model.
  \begin{itemize}
    \item 0: Turn the logistic terms off.
    \item 1: Turn the logistic terms on. 
      This will change %the $G$-equation of the OPN model to (see the next section):
      %\begin{displaymath}
      %  \pp{(r^2G)}{t} + \po{r}\left[r^2 G V\right] = {\color{blue}\lambda}\left(\frac{r}{R(t)}\right)^{\color{blue}q}\left[1+\frac{{\color{blue}b'}P_2}{{\color{blue}b}+P_2}\right]r^2G - {\color{blue}\mu}\left[1+\frac{{\color{blue}a'}P_1}{{\color{blue}a}+P_1}\right]r^2G\;;
      %\end{displaymath}
      %and 
      the $V$-equation to:
      \begin{displaymath}
        \frac{1}{r^2}\po{r}\left[r^2\left(V+{\color{blue}\alpha} M\pp{A}{r}\right)\right] = {\color{blue}\lambda}\left(\frac{r}{R(t)}\right)^{\color{blue}q}\left[1+\frac{{\color{blue}b'}P_2}{{\color{blue}b}+P_2}\right] G(1-G) - {\color{blue}\delta_1} H - {\color{blue}\delta_2} M\;.
      \end{displaymath}
  \end{itemize}
\end{itemize}
%For the second point, if a uniform dissipation rate for $N$ is to be used, you can delete the lines for \texttt{n\_star}, \texttt{rho\_mnus}, and \texttt{rho\_plus}, and use the following line instead:
%\begin{lstlisting}[backgroundcolor = \color{lightgray}, language = bash, xleftmargin = 0.1cm, framexleftmargin = 0.1cm]
%rho = 0.8
%\end{lstlisting}
%Here 0.8 is to be replaced by any value assigned to $\rho$.
The rest of the first section prescribes model parameters, see Section~\ref{sec:model} for more details.
If a variable is not assigned in the input file, the default value will be used; to get an idea what the default value is for these variables, take a look at the first function in \texttt{\$TDHOME/TDSolve/src/TDInput.cc}.

The second section features simulation parameters, and they're explained below:
\begin{itemize}
\item \texttt{ic}: The initial condition tag. 
It refers to a tag that is hard-coded in the program. 
For example, the option 0 of the problem 0 corresponds to the following initial data:
\begin{align*}
&N(r,0) = N_{\textrm{bc}}\;,\ G(r,0) = 0.84\;,\ H(r,0) = 0.16-N_{\textrm{bc}}\;; \\
&A(r,0) = 10^3\exp(-r^2)\;.
\end{align*}
For more details, refer to the function:
\begin{lstlisting}[backgroundcolor = \color{lightgray}, language = bash, xleftmargin = 0.1cm, framexleftmargin = 0.1cm]
void TDSolve::set_ic(void)
\end{lstlisting}
in \texttt{\$TDHOME/TDSolve/src/TDSolveFVM.cc}.
This is also where you should modify to introduce more initial conditions (and re-compile the program).
\item \texttt{r\_0}: The initial radius of the tumor.
\item \texttt{Nx}: The number of uniform cells to discretize the domain $r\in[0, R(t)]$.
\item \texttt{A\_Multiplier}: The variable $A$ is solved on a larger domain $r\in[0, M_AR(t)]$, and the multiplier is given by $M_A$.
\item \texttt{cfl}: The Courant number $\alpha_{\textrm{cfl}}$ to compute the time step size. For explicit method its value should be no more than $1.0$ and the suggested value is $0.4$; for implicit method its value can be increased to $\sim 20.0$.
\item \texttt{Tinit}: The initial time -- the default is zero.
\item \texttt{Tterm}: The terminal time -- the difference between the terminal and the initial times is the time span of the numerical simulations.
\item \texttt{Rterm}: The terminal radius -- the computation will stop when Rterm is reached.
\item \texttt{Rdead}: The diminishing radius -- the computation will stop when Rdead is reached. It is required that \texttt{Rterm} $>$ \texttt{R\_0} $>$ \texttt{Rdead}.
\item \texttt{Tfreq}: The frequency at which the solution is sampled and written to the solution files.
\item \texttt{time\_method}: Time-integrator; 0 for a first-order method and 1 for a second-order method.
\item \texttt{max\_iter}: Maximum allowed number of iterations. The program will end whenever this number is reached or the terminal time is achieved, whichever is earlier.
\item \texttt{dt\_init}: The beginning time step size $dt_{\textrm{init}}$ (see below).
\item \texttt{dt\_growth}: The growth rate $\alpha_{\textrm{growth}}$ of time step size.
\end{itemize}
Among the list above, three factors are used to determine the time step size.
Let the time-step size determined from the Courant number be $dt^k_{\textrm{cfl}}$, i.e., 
\begin{displaymath}
dt^k_{\textrm{cfl}} = \min_i\frac{\alpha_{\textrm{cfl}}V^k_i}{h}\;,
\end{displaymath}
$h$ being the cell size in non-dimensional coordinate.
Then at the firs step, the actual time-step size $dt^0$ is:
\begin{displaymath}
dt^0 = \min(dt_{\textrm{init}}\;,\ dt^0_{\textrm{cfl}})\;;
\end{displaymath}
in later time steps, we have:
\begin{displaymath}
dt^k = \min(\alpha_{\textrm{growth}}dt^{k-1}\;,\ dt^k_{\textrm{cfl}})\;.
\end{displaymath}
This technique is known as {\it ramp-up}.

The third section includes the options for numerical methods:
\begin{itemize}
\item \texttt{framework}: The general framework of the numerical method.
  \begin{itemize}
    \item Set it to 0 to use a finite-difference method (cannot capture shocks, not updated for recent models anymore).
    \item Set it to 1 for finite-volume method (always use this one).
  \end{itemize}
\item \texttt{space\_cell}: The numerical method to update the cell numbers. 
  \begin{itemize}
    \item Set it to 0 for an explicit first-order upwind method.
    \item Set it to 1 for an implicit first-order upwind method.
    \item Set it to 2 for an explicit second-order MUSCL method.
  \end{itemize}
\item \texttt{space\_chem}: The numerical method to update the chemotaxis concentration. 
  \begin{itemize}
    \item Set it to 0 for an implicit method: first-order upwind for advection and second-order central difference for diffusion.
    \item Set it to 1 for an IMEX method: explicit second-order MUSCL for advection and implicit second-order central difference for diffusion.
  \end{itemize}
\item \texttt{vel}: The numerical method to compute the velocity.
Only one option 0 is available right now -- a second-order integral form is used to compute the tumor velocity.
\item \texttt{lsysType}: The linear solver option.
Many methods are supported, including direct solver (LU factorization) and Krylov sub-space methods.
For this problem, set it to \texttt{DIR} is good enough.
\item \texttt{lsysPrint}: The option 0 turn off the diagnostic information of linear solve; the option 1 turn it on.
\end{itemize}

The last section specifies the output files:
\begin{itemize}
\item \texttt{outputDataFile}: The file for all the solutions at all in-tumor cells; the solutions are sampled according to \texttt{Tfreq}.
\item \texttt{radiusDataFile}: The file for time histories of radius, tumor boundary velocity, and the chemotaxis concentration gradient at the tumor boundary.
These information is output every time step.
\end{itemize}
More about the output files are given in the next section.

\subsection{The Output Files}
Running the program will generate two files, one is solution file and the other one is time history file.
The solution file (\texttt{sol.dat} in the previous example) starts with a line describing the variables in each subsequent section. 
The rest of the file is sub-divided into $N_t$ parts, every single part corresponds to the solutions at a specific output time.
Each part start with a line containing two double variables and two integers, such as:
\begin{lstlisting}[backgroundcolor = \color{lightgray}, language = bash, xleftmargin = 0.1cm, framexleftmargin = 0.1cm]
0.02000000000000000042	0.00438871438561059483	50	51
\end{lstlisting}
These numbers are: the solution time, the time step size, the number of cells, and the number of cell boundaries; the difference between the latter two is always one.
This line is followed by, in this examples, 50 lines describing the cell numbers and the chemotaxis concentration at each cell; we present an example with less precision to save space:
\begin{lstlisting}[backgroundcolor = \color{lightgray}, language = bash, xleftmargin = 0.1cm, framexleftmargin = 0.1cm]
0.00202 0.55136 0.10392 1.08261 212.88046
\end{lstlisting}
From the left to the right, these numbers are: the dimensional cell center coordinate, $G$, $H$, $N$, and $A$.
The output variables differ among the problems, which are determined by the input variable \texttt{prob\_id}:
\begin{itemize}
\item \texttt{prob\_id}=0 (PDGF): $G$, $H$, $N$, and $A$.
\item \texttt{prob\_id}=1 (OPN): $G$, $H$, $M$, $A$, $P_1$, and $P_2$.
\end{itemize}
This part ends with 51 addition lines describing the velocity at each cell faces; each line contains two numbers, the dimensional cell-face coordinate and the velocity at this point.

The time history file (\texttt{rad.dat} in the previous example) is much simpler, it contains $T$ lines with $T$ being the total number of time steps in the simulation.
Each line contains four numbers:
{\tiny
\begin{lstlisting}[backgroundcolor = \color{lightgray}, language = bash, xleftmargin = 0.1cm, framexleftmargin = 0.1cm]
0.00000100000000000000	0.20000005157812633660	0.05194605632028042458	396.22296116163556689571
\end{lstlisting}
}\noindent
From the left to the right, they are: solution time $t$, radius $R(t)$, velocity at tumor boundary $V(R(t),t)$, and the negative chemotaxis gradient at the boundary $-\pp{A}{r}(R(t),t)$.

\subsection{Reading the Solution Files using \texttt{Matlab}}
Clearly, \texttt{Matlab} (or \texttt{Octave}) is required for this part.
A few scripts are in the folder \texttt{\$TDHOME/tools}; since now we only use the finite volume method to solve the problem, ignore all scripts without \texttt{\_fvm\_} in the file name.

Copy the scripts to the directory where \texttt{Matlab} is started; alternatively, you can add the folder \texttt{\$TDHOME/tools} to the path of \texttt{Matlab}.
Among all, the most important script is \texttt{read\_data\_fvm\_1d\_tumordyn.m}; it reads the problem identifier and the solution file (\texttt{sol.dat}) and return the solutions in \texttt{Matlab} format.
In \texttt{Matlab}:
\begin{lstlisting}[backgroundcolor = \color{lightgray}, language = Matlab, xleftmargin = 0.1cm, framexleftmargin = 0.1cm]
> [d_c,d_f,time] = read_data_fvm_1d_tumordyn(0,'sol.dat');
\end{lstlisting}
will create three arrays:
\begin{itemize}
\item \texttt{d\_c}: An $N_t\times N_d\times N_x$ array, where $N_t$ is the number of output times, $N_d$ corresponds to the numbers described before ($5$ for the PDGF model and $7$ for the OPN model), and $N_x$ is the number of cells (50 in the previous example).
\item \texttt{d\_f}: An $N_t\times2\times(N_x+1)$ array, where $N_t$ is the number of output times, $2$ corresponds to the cell-face coordinate and velocity, and $N_x$ is the number of cells (50 in the previous example).
\item \texttt{time}: An $1\times N_t$ array storing the actual solution times.
\end{itemize}
The scripts \texttt{plot\_data\_fvm\_1d\_tumordyn.m} and \texttt{make\_mov\_data\_fvm\_1d\_tumordyn.m} plot and make a movie of the solution file, respectively; they are not quite used now.
The script \texttt{proc\_voldata\_fvm\_1d\_tumordyn.m}, however, post-process the solution file and create global data such as total number of $G$, $H$, and $N$ in the tumor, their sum, and their proportions at each data output time.
Particularly, calling the following function in \texttt{Matlab}:
\begin{lstlisting}[backgroundcolor = \color{lightgray}, language = Matlab, xleftmargin = 0.1cm, framexleftmargin = 0.1cm]
> data = proc_voldata_fvm_1d_tumordyn(0,'sol.dat');
\end{lstlisting}
will create an $N_t\times8$ array \texttt{data}, with the $8$ columns being: 
\begin{itemize}
  \item PDGF model: (1) solution time, (2) total $G$, (3) total $H$, (4) total $N$, (5) total $G+H+N$, (6) portion of $G$, (7) portion of $H$, and (8) portion of $N$.
  \item OPN model: (1) solution time, (2) total $G$, (3) total $H$, (4) total $M$, (5) total $G+H+M$, (6) portion of $G$, (7) portion of $H$, and (8) portion of $M$.
\end{itemize}

\section{Models and Their Parameter Lists}
\label{sec:model}

\subsection{Model Identifiers and Governing Equations}
The models are determined by the variable \texttt{prob\_id}; and each model contains parameters that are marked in {\color{blue} blue} fonts:
\begin{itemize}
\item \texttt{prob\_id}=0: PDGF model (default):
  \begin{align*}
    &\pp{(r^2G)}{t} + \po{r}\left[r^2 G V\right] = {\color{blue}\lambda}\left(\frac{r}{R(t)}\right)^{\color{blue}q} r^2G - {\color{blue}\mu} r^2G\;, \\
    &\pp{(r^2H)}{t} + \po{r}\left[r^2H V\right] = {\color{blue}\mu} r^2G - {\color{blue}\delta} r^2H\;, \\
    &\pp{(r^2N)}{t} + \po{r}\left[r^2N\left(V+{\color{blue}\alpha}\pp{A}{r}\right)\right] = - {\color{blue}\rho} r^2N\;,\\
    &\pp{(r^2A)}{t} = {\color{blue}D}\po{r}\left(r^2\pp{A}{r}\right) + \frac{\chi_{0\le r\le R}{\color{blue}m}r^2G}{{\color{blue}\beta}/{\color{blue}\theta}+G} - {\color{blue}\gamma} r^2A\;,\\
    &\frac{1}{r^2}\po{r}\left[r^2\left(V+{\color{blue}\alpha} N\pp{A}{r}\right)\right] = {\color{blue}\lambda}\left(\frac{r}{R(t)}\right)^{\color{blue}q} G - {\color{blue}\delta} H - {\color{blue}\rho} N\;.
  \end{align*}
\item \texttt{prob\_id}=1: OPN model:
  \begin{align*}
    &\pp{(r^2G)}{t} + \po{r}\left[r^2 G V\right] = {\color{blue}\lambda}\left(\frac{r}{R(t)}\right)^{\color{blue}q}\left[1+\frac{{\color{blue}b'}P_2}{{\color{blue}b}+P_2}\right]r^2G - {\color{blue}\mu}\left[1+\frac{{\color{blue}a'}P_1}{{\color{blue}a}+P_1}\right]r^2G - {\color{blue} \omega}r^2GM\;, \\
    &\pp{(r^2H)}{t} + \po{r}\left[r^2H V\right] = {\color{blue}\mu}\left[1+\frac{{\color{blue}a'}P_1}{{\color{blue}a}+P_1}\right]r^2G - {\color{blue}\delta_1} r^2H\;, \\
    &\pp{(r^2M)}{t} + \po{r}\left[r^2M\left(V+{\color{blue}\alpha}\pp{A}{r}\right)\right] = - {\color{blue}\delta_2} r^2M\;,\\
    &\pp{(r^2A)}{t} = {\color{blue}D_3}\po{r}\left(r^2\pp{A}{r}\right) + \frac{\chi_{0\le r\le R}{\color{blue}\alpha_3}r^2G}{{\color{blue}\beta_3}/{\color{blue}\theta}+G} - {\color{blue}\gamma_3} r^2A\;,\\
    &\pp{(r^2P_1)}{t} = {\color{blue}D_1}\po{r}\left(r^2\pp{P_1}{r}\right) + \frac{{\color{blue}\alpha_1}r^2M}{{\color{blue}\beta_1}/{\color{blue}\theta}+M} - {\color{blue}\gamma_1} r^2P_1\;,\\
    &\pp{(r^2P_2)}{t} = {\color{blue}D_2}\po{r}\left(r^2\pp{P_2}{r}\right) + \frac{{\color{blue}\alpha_2}r^2G}{{\color{blue}\beta_2}/{\color{blue}\theta}+G} - {\color{blue}\gamma_2} r^2P_2\;,\\
    &\frac{1}{r^2}\po{r}\left[r^2\left(V+{\color{blue}\alpha} M\pp{A}{r}\right)\right] = {\color{blue}\lambda}\left(\frac{r}{R(t)}\right)^{\color{blue}q}\left[1+\frac{{\color{blue}b'}P_2}{{\color{blue}b}+P_2}\right] G - {\color{blue}\delta_1} H - {\color{blue}\delta_2} M\;.
  \end{align*}
\end{itemize}

\subsection{Model Parameters and Input Files}
The parameters differ among models given before. 
In the next table (Table~\ref{tb:par}) we summarize the correspondence between the parameter described in the equations, their input file variables, the \texttt{C} variables, and their default values.
\begin{table}[ht]\centering
\setlength{\tabcolsep}{3pt}
\caption{Parameter summary.}
\label{tb:par}
\begin{tabular}{|l||c|c|l|c|}
\hline
Input file        & PDGF      & OPN        & \texttt{C} variable & Default value \\ \hline \hline
\texttt{lambda}   & $\lambda$ & $\lambda$  & \texttt{lambda}     & $1.0$         \\ \hline
\texttt{mu}       & $\mu$     & $\mu$      & \texttt{mu}         & $1.0$         \\ \hline
\texttt{delta\_h} & $\delta$  & $\delta_1$ & \texttt{delta\_h}   & $1.0$         \\
\texttt{delta}    &           &            &                     &               \\
\texttt{delta\_1} &           &            &                     &               \\ \hline
\texttt{alpha}    & $\alpha$  & $\alpha$   & \texttt{alpha}      & $1.0$         \\ \hline
\texttt{theta}    & $\theta$  & $\theta$   & \texttt{theta}      & $1.0$         \\ \hline
\texttt{delta\_n} & $\rho$    & $\delta_2$ & \texttt{delta\_n}   & $0.0$         \\
\texttt{rho}      &           &            &                     &               \\
\texttt{delta\_2} &           &            &                     &               \\ \hline
\texttt{D}        & $D$       & $D_3$      & \texttt{D}          & $1.0$         \\
\texttt{D\_3}     &           &            &                     &               \\ \hline
\texttt{alpha\_a} & $m$       & $\alpha_3$ & \texttt{alpha\_a}   & $1.0$         \\
\texttt{m}        &           &            &                     &               \\
\texttt{alpha\_3} &           &            &                     &               \\ \hline
\texttt{beta\_a}  & $\beta$   & $\beta_3$  & \texttt{beta\_a}    & $1.0$         \\
\texttt{beta}     &           &            &                     &               \\
\texttt{beta\_3}  &           &            &                     &               \\ \hline
\texttt{gamma\_a} & $\gamma$  & $\gamma_3$ & \texttt{gamma\_a}   & $1.0$         \\
\texttt{gamma}    &           &            &                     &               \\
\texttt{gamma\_3} &           &            &                     &               \\ \hline
\texttt{D\_1}     &           & $D_1$      & \texttt{D\_1}       & $0.01$        \\ \hline
\texttt{alpha\_1} &           & $\alpha_1$ & \texttt{alpha\_1}   & $10.0$        \\ \hline
\texttt{beta\_1}  &           & $\beta_1$  & \texttt{beta\_1}    & $1000.0$      \\ \hline
\texttt{gamma\_1} &           & $\gamma_1$ & \texttt{gamma\_1}   & $1.0$         \\ \hline
\texttt{D\_2}     &           & $D_2$      & \texttt{D\_2}       & $0.01$        \\ \hline
\texttt{alpha\_2} &           & $\alpha_2$ & \texttt{alpha\_2}   & $10.0$        \\ \hline
\texttt{beta\_2}  &           & $\beta_2$  & \texttt{beta\_2}    & $1000.0$      \\ \hline
\texttt{gamma\_2} &           & $\gamma_2$ & \texttt{gamma\_2}   & $1.0$         \\ \hline
\texttt{a}        &           & $a$        & \texttt{a}          & $100.0$       \\ \hline
\texttt{b}        &           & $b$        & \texttt{b}          & $100.0$       \\ \hline
\texttt{a\_prime} &           & $a'$       & \texttt{a\_prime}   & $1.0$         \\ \hline
\texttt{b\_prime} &           & $b'$       & \texttt{b\_prime}   & $1.0$         \\ \hline
\texttt{q}        & $q$       & $q$        & \texttt{q}          & $0.0$         \\ \hline
\texttt{omega}    &           & $\omega$   & \texttt{omega}      & $0.0$         \\ \hline
\end{tabular}
\setlength{\tabcolsep}{1pt}
\end{table}

\section{Case Study: A Parametric Study of the PDGF Model}
Lastly, the folder \texttt{\$TDHOME/manual/demo} contains a package \texttt{par\_m\_pdgf.tar.gz}.
This is an example of a small project of parametric study of the tumor growth model with different values of $m$ (say $1e+5$, $2e+5$, $4e+5$, and $8e+5$), keeping all the other parameters fixed.
A similar package for the OPN model is given by \texttt{par\_alpha\_opn.tar.gz}; using this package follows the same steps.

To begin with, let us copy this package to \texttt{\$TDHOME/TDRun} and expand it:
\begin{lstlisting}[backgroundcolor = \color{lightgray}, language = bash, xleftmargin = 0.1cm, framexleftmargin = 0.0cm]
$ cd $TDHOME/TDRun
$ cp $TDHOME/TDSolve/manual/demo/par_m.tar.gz ./
$ tar xvf par_m.tar.gz
\end{lstlisting}
This will create a folder \texttt{TestProjPDGF} that contains a input file, a number of shell scripts, and a folder \texttt{data/}; this latter folder contain a few \texttt{Matlab} scripts to make plots.

\subsection{Run the Simulations One-By-One}
The first method to run the set of simulations with different $m$ values is to use the script \texttt{runall.sh} as follows: While in \texttt{\$TDHOME/TDRun/}:
\begin{lstlisting}[backgroundcolor = \color{lightgray}, language = bash, xleftmargin = 0.1cm, framexleftmargin = 0.0cm]
$ cd TestProjPDGF
$ sh runall.sh
\end{lstlisting}
This script will loop over all the four $m$ values, replace the string \texttt{VALM} in the template input file \texttt{input\_raw.hcl} by these values, run the simulation, and rename and copy the results to the sub-directory \texttt{data/}.
All the four simulations are done one after another, so only one processor is utilized.
Note that the shell script \texttt{runall.sh} uses the GNU stream editor \texttt{sed}, which is mostly likely already available on the system.

\subsection{Run the Simulations in Parallel}
The solver itself is not parallelized; however, if more the one simulations need to be done, there is no reason why not to make use of the multiple processors that are available on the machine.
A simple tool to achieve this is to use GNU Parallel ({\color{blue} \url{https://www.gnu.org/software/parallel/}}), it is available in the package management tool for your system.
With this tool installed, follow the next commands to run the simulations: While in \texttt{\$TDHOME/TDRun/}:
\begin{lstlisting}[backgroundcolor = \color{lightgray}, language = bash, xleftmargin = 0.1cm, framexleftmargin = 0.0cm]
$ cd TestProjPDGF
$ parallel sh run_par.sh ::: 1e+5 2e+5 4e+5 8e+5
\end{lstlisting}
Then all the free processors will be utilized to run these simulations.

\subsection{Post-process the Data}
Using either method, four sets of solutions are created, properly named, and copied to the folder \texttt{data/}.
Go into this directory and start \texttt{Matlab}, which will load the \texttt{startup.m} and add to the path the \texttt{tools/} discussed before.
If you're using \texttt{Octave} instead, comment out the line \texttt{beep off} in \texttt{startup.m}.
Then in \texttt{Matlab} or \texttt{Octave} interface:
\begin{lstlisting}[backgroundcolor = \color{lightgray}, language = Matlab, xleftmargin = 0.1cm, framexleftmargin = 0.1cm]
> plot_rads;
\end{lstlisting}
will plot the radius growth histories for the set of simulations; and
\begin{lstlisting}[backgroundcolor = \color{lightgray}, language = Matlab, xleftmargin = 0.1cm, framexleftmargin = 0.1cm]
> plot_n_port;
\end{lstlisting}
will plot the portion of the immune cells ($N$) as a function of time.
Both scripts also save the plotted figures into \texttt{eps} format.

\subsection{Final Remark}
In this project, the scripts \texttt{run\_1.sh}, \texttt{run\_2.sh}, and \texttt{startup.m} in the sub-directory \texttt{data/} assume that you're following exactly this manual, particularly the directory hierarchy.
If this is not the case, please modify these three scripts accordingly.

\end{document}

